% \chapter{Diskussion}

\chapter{Zusammenfassung und Ausblick}
\label{sec:za}
In dieser Bachelorarbeit wurde die generation von THz-Strahlung durch zwei organische Kristalle geprüft
und ob diese ZnTe als Generationskristall ersetzen können. 
Der zugrunde liegende Aufbau ist ein THz-TDS-AUbau mit einer EOS-Detektion. 
Zu erst wurde die Funktionalität des Lock-In-Verstärkers getestet bei welchem das Verhalten 
des Rauschens auf verschiedene zeitkonstanten $\tau$ geprüft wurde. 
Dazu wurde der einfluss von verschiedenen Stromversorgungen der Photodioden geprüft. 
Es ergibt sich, dass die Batterien hingegen der Erwartung, mehr Rauschen als die Netzspannung erzeugen. 
Als Ursache dafür ist es möglich, dass der Batterie Aufbau, welcher nur ein Prototyp ist nicht vollständig funktioniert. 
Für den weiteren Aufbau wurde die genaue Strahlgröße mithilfe des Rasierklingen-Verfahrens bestimmt und die Leistungsdichte des Strahls bestimmt.
Für den ZnTe-Kristall wird wie auch für den OH1 Kristall die Lineare Abhängigkeit von der Pump-Leistung gezeigt. 
Des Weiteren besitzt der THz-Puls erzeugt von dem ZnTe eine Bandbreite von \SI{2.4}{THz}, wobei 
Zwei Absorptionslinien im Frequenzspektrum bei \SI{1.2}{THz} und \SI{1.7}{THz} sichtbar sind. 
Als einschränkender Faktor sei erwähnt, dass die Messwerte zwischen \SIrange{30}{4}{mW} verzerrt sind. 
 
Ab einer Pumpleistung von \SI{7.5}{mW} ist ein THz-Puls erzeugt von dem OH1-Kristall Messbar. 
Der Kristall erzeugt also auch bei einer Pump-Wellenlänge von \SI{808}{nm} THz-Strahlung. 
Es ist zu erwarten gewesen, dass der Anstieg der Maximalen Amplitude zur angelegten Pumpleistung höher ist als bei ZnTe. 
Das dies nicht der Fall ist lässt sich durch den Detektionskristall erklären, da dieser ein ZnTe Kristall ist und somit die 
vermutlich größere Bandbreite, welchhe vom OH1-Kristall erzeugt wurde, nicht in ein Messbares Signal umwandeln kann. 
Somit lässt sich darüber erkläen warum die beiden gerade nahe bei einander sind. 

Der BNA-Kristall hat kein Signal über dem Rauschhniveau erzeugt bei bis zu \SI{480}{mW} Pumpleistung. 
Als Ursache für dies ist die Probenqualität anzuführen, da diese nur eine polykristline Pulverprobe war und innerhalb der Probe kleine 
einzel Kristalle vorhanden sind. 
Dadurch erzeugen die einzel Kristalle destruktive Interverenz durch den ungeraden Tensor-Rang von $\chi^{(2)}$. 
Es ist zu betonen, dass das über das Material generell keine Aussage getroffen werden kann, sondern die Proben größe und Qualität Ursache 
für das nicht erscheinen eines Signals ist.

Abschließend lässt sich zu der Nutzbarkeit sagen, dass über BNA bei der Derzeitigen Probenqualität keine Aussage über die THz-Generation 
getroffen werden kann. 
Es ist eine größere und zusammenhängende Kristallfäche von nöten um mehr über BNA sagen zu können. 
ZnTe bleibt bei dem Aufbau der effektivste Emitter, dies ist aber sehr beeinflusst durch den Detektiosnkristall welcher auch ein ZnTe-Kristall ist. 
Hierbei kann in zukunft auf andere Kristalle wie GaSe ausgewichen werden, welche eine höhere Bandbreiten Detektion besitzen. 
Dadurch kann das möglicherweise größere Frequenzspektrum des erzeugten THz-Pulses durch den OH1-Kristall vollständig detektiert werden. 
Des Weiiteren ist auch die Probengröße des OH1-Kristalls nicht zufriedenstellend. 
Eine größere Optischaktive Fläche um die gesamte Pumpleistung nutzen ist für die Zukuft anstrebsam. 
Hinzu kommt, dass OH1 nicht für eine Pumpwellenlänge von \SI{808}{nm} Optimiert ist. 
Die abhängigkeit von der erzeugten THz-Feldstärke von der Pumpwellenlänge könnte in zukunft untersucht werden. 
Abseits der Proben ist die Stromversorgung der Photodioden auch ein ansatzpunkt, da gezeigt wurde das die genuzten Batterien keinen Mehrwert für 
die Rauschspannung bieten. 
Hier kann in der zusammenarbeit mit der Elektronischenwekstatt nach alternativen gesucht werden. 
Ein anderer Punkt sind die gemessenen Wasserabsorptionslinien, welche im Frequenzspektrum von ZnTe klar zu sehen sind.
Um diese zu enfernen ist es von Vorteil den bereich des Aufbaus in welchem THz-Strahlung propagiert mit einer Purgingbox zu versehen. 
In diese Purgingbox kann Trockenluft oder Stickstoff eingefürt werden, 
der Wasserdamp die THz-Strahlung nichht absorbiert. 
